
\documentclass{article}
\usepackage{colortbl}
\usepackage{makecell}
\usepackage{multirow}
\usepackage{supertabular}

\begin{document}

\newcounter{utterance}

\centering \large Interaction Transcript for game `eqbench', experiment `eq\_bench', episode 132 with reflexion{-}Qwen3.5{-}27B.
\vspace{24pt}

{ \footnotesize  \setcounter{utterance}{1}
\setlength{\tabcolsep}{0pt}
\begin{supertabular}{c@{$\;$}|p{.15\linewidth}@{}p{.15\linewidth}p{.15\linewidth}p{.15\linewidth}p{.15\linewidth}p{.15\linewidth}}
    \# & \multicolumn{2}{c}{Player} && \multicolumn{2}{c}{Game Master} \\
    \hline

    \theutterance \stepcounter{utterance}  
    & & & \multicolumn{4}{p{0.6\linewidth}}{
        \cellcolor[rgb]{0.9,0.9,0.9}{
            \makecell[{{p{\linewidth}}}]{
                \texttt{\tiny{[P1$\langle$GM]}}
                \texttt{You have attempted this same task before and did not succeed. These are the plans you wrote after those attempts.} \\
\texttt{Trial \#1: New plan: My previous attempt failed because I misjudged the primary emotional driver of Arthur's final state. I scored "Frustrated" highest (8) and "Misunderstood" second (6), but I underestimated "Hopeful" (2). Arthur's dialogue is a persistent, escalating attempt to repair the relationship ("I'm trying to change," "I can try to make it up to you"). Despite Sam's harsh rejection, Arthur's final stance is one of desperate optimism that his apology will eventually be accepted, not just frustration at the rejection. Next time, I will prioritize the character's *intent* and *goal* (reconciliation) over the immediate *reaction* to the conflict. I will assign the highest score to "Hopeful" to reflect his persistent belief that he can fix things, lower "Frustrated" to reflect that he is accepting the rejection rather than lashing out, and adjust "Misunderstood" to reflect that he acknowledges his fault ("I messed up") rather than feeling unfairly judged.} \\
\\ 
\texttt{Trial \#2: New plan: My previous attempt failed because I overestimated "Hopeful" (8) and underestimated "Frustrated" (6), likely misinterpreting Arthur's persistence as optimism rather than desperation. Arthur's final lines ("I can't take it back," "I can try") show he is hitting a wall; Sam's rhetorical questions ("How, Arthur?") explicitly shut down his attempts, leaving him with no path forward. This creates high frustration, not hope. Additionally, Arthur admits his faults ("I messed up"), so "Misunderstood" should be very low. Next time, I will prioritize the *outcome* of the interaction: since Arthur's efforts are explicitly rejected and he has no solution to offer, "Frustrated" must be the highest score. "Hopeful" should be low because the dialogue ends with a dead end, not a breakthrough. "Misunderstood" should be near zero as he accepts blame, and "Angry" should remain low as he is not lashing out. I will assign scores reflecting this: Frustrated (9), Hopeful (2), Misunderstood (1), Angry (1).} \\
\\ 
\texttt{Trial \#3: New plan: My previous attempt failed because I misinterpreted Arthur's persistence as "Frustrated" (scoring it 9) rather than recognizing it as a desperate, lingering "Hopeful" attempt to reconcile. While Sam is rejecting him, Arthur's final lines ("I can try to make it up to you") are not expressions of hitting a wall, but rather a continued, albeit weak, offer of repair. The core of his emotional state is the hope that this specific attempt might finally work, despite the odds. "Frustrated" is present but secondary to his driving motivation. "Misunderstood" is low because he admits fault, and "Angry" is negligible. Next time, I will prioritize the character's *active goal* (reconciliation) over the *passive reaction* to rejection. I will assign the highest score to "Hopeful" (e.g., 8) to reflect his persistent belief in fixing the relationship, lower "Frustrated" (e.g., 4) to acknowledge the difficulty without making it the dominant emotion, and keep "Misunderstood" and "Angry" very low (0{-}1).} \\
\\ 
\\ 
\\ 
\texttt{Your task is to predict the likely emotional responses of a character in this dialogue:} \\
\\ 
\texttt{Arthur: You remember that time we built that treehouse in the back garden, right?} \\
\texttt{Sam: What's that got to do with anything?} \\
\texttt{Arthur: I'm just saying, we were a team, you and me. We built something out of nothing. } \\
\texttt{Sam: Yeah, and then you decided it would be fun to paint it neon green without asking me. } \\
\texttt{Arthur: It was a creative decision, Sammy. } \\
\texttt{Sam: It was a stupid decision, Arthur. Just like your speech at my wedding.} \\
\texttt{Arthur: That's not fair, Sam. I messed up, I know. But I'm trying to apologize here.} \\
\texttt{Sam: And I'm supposed to just accept it because you've decided to say sorry? } \\
\texttt{Arthur: No, not because I've decided to. Because I mean it, Sam. I'm sorry.} \\
\texttt{Sam: You're always sorry, Arthur. After every stupid, thoughtless thing you do. } \\
\texttt{Arthur: I know, and I'm trying to change that. I'm trying to be better.} \\
\texttt{Sam: Better? You think saying sorry makes you better? You ruined the most important day of my life, Arthur.} \\
\texttt{Arthur: I know, and I can't take it back. But I can try to make it up to you, Sam.} \\
\texttt{Sam: How, Arthur? By bringing up old memories? By reminding me of when we were kids?} \\
\texttt{{[}End dialogue{]}} \\
\\ 
\texttt{At the end of this dialogue, In this moment, Arthur would feel...} \\
\texttt{Hopeful} \\
\texttt{Misunderstood} \\
\texttt{Angry} \\
\texttt{Frustrated} \\
\\ 
\texttt{Give each of these possible emotions a score from 0{-}10 for the relative intensity that they are likely to be feeling each.} \\
\\ 
\texttt{You must output in the following format, including headings (of course, you should give your own scores), with no additional commentary:} \\
\\ 
\texttt{Hopeful: <score>} \\
\texttt{Misunderstood: <score>} \\
\texttt{Angry: <score>} \\
\texttt{Frustrated: <score>} \\
\\ 
\\ 
\texttt{{[}End of answer{]}} \\
\\ 
\texttt{Remember: zero is a valid score, meaning they are likely not feeling that emotion. You must score at least one emotion > 0.} \\
\\ 
\texttt{Your answer:} \\
            }
        }
    }
     \\ \\

    \theutterance \stepcounter{utterance}  
    & \multicolumn{4}{p{0.6\linewidth}}{
        \cellcolor[rgb]{0.561,0.737,0.561}{
            \makecell[{{p{\linewidth}}}]{
                \texttt{\tiny{[P1$\rangle$GM]}}
                \texttt{Hopeful: 2} \\
\texttt{Misunderstood: 1} \\
\texttt{Angry: 0} \\
\texttt{Frustrated: 9} \\
            }
        }
    }
    & & \\ \\

    \theutterance \stepcounter{utterance}  
    & & & \multicolumn{2}{p{0.3\linewidth}}{
        \cellcolor[rgb]{0.9,0.9,0.9}{
            \makecell[{{p{\linewidth}}}]{
                \texttt{\tiny{[GM$|$GM]}}
                \texttt{Hopeful: 2} \\
\texttt{Misunderstood: 1} \\
\texttt{Angry: 0} \\
\texttt{Frustrated: 9} \\
            }
        }
    }
    & & \\ \\

    \theutterance \stepcounter{utterance}  
    & & & \multicolumn{2}{p{0.3\linewidth}}{
        \cellcolor[rgb]{0.9,0.9,0.9}{
            \makecell[{{p{\linewidth}}}]{
                \texttt{\tiny{[GM$|$GM]}}
                \texttt{Hopeful: 3} \\
\texttt{Misunderstood: 7} \\
\texttt{Angry: 2} \\
\texttt{Frustrated: 8} \\
            }
        }
    }
    & & \\ \\

    \theutterance \stepcounter{utterance}  
    & & & \multicolumn{2}{p{0.3\linewidth}}{
        \cellcolor[rgb]{0.9,0.9,0.9}{
            \makecell[{{p{\linewidth}}}]{
                \texttt{\tiny{[GM$|$GM]}}
                \texttt{game\_result = LOSE} \\
            }
        }
    }
    & & \\ \\

    \theutterance \stepcounter{utterance}  
    & & \multicolumn{4}{p{0.6\linewidth}}{
        \cellcolor[rgb]{0.95,0.95,0.95}{
            \makecell[{{p{\linewidth}}}]{
                \texttt{\tiny{[P1$\leftrightarrow$P]}}
                \texttt{Post{-}episode Reflexion (not sent to GM):} \\
\texttt{New plan: My previous attempts failed because I oscillated between interpreting Arthur's persistence as "Hopeful" (Trial 1 \& 3) and "Frustrated" (Trial 2), failing to recognize the nuance of his specific emotional state in the final moment. Arthur is not "Hopeful" because Sam's final rhetorical questions ("How, Arthur?") explicitly dismantle his offers, leaving no room for optimism. He is not primarily "Frustrated" in an angry sense, as he remains calm and apologetic. The dominant emotion is "Misunderstood" (or rather, the feeling that his *intent* to change is not being seen or accepted, even if he admits his *actions* were wrong). He feels his sincere effort ("I mean it") is being dismissed as just another empty apology ("You're always sorry"). Therefore, I will assign the highest score to "Misunderstood" (8) to reflect the gap between his internal sincerity and Sam's external rejection. "Frustrated" will be secondary (5) due to the inability to bridge this gap. "Hopeful" will be low (2) as the dialogue ends on a dead end, and "Angry" will be 0.} \\
            }
        }
    }
    & \\ \\

\end{supertabular}
}

\end{document}
